\begin{figure}[tb]
\centering
\begin{tikzpicture}[>=Latex,font=\small,scale=.95]
\draw[->] (-3.3,0)--(3.4,0) node[right]{\(i_d\)};
\draw[->] (0,-2.5)--(0,2.7) node[above]{\(i_q\)};
\foreach \c in {.45,.85,1.25,1.65}
 {
  \draw[Purple!65,domain=.55:3,samples=70] plot ({\x},{\c/\x});
  \draw[Purple!35,domain=-3:-.55,samples=70] plot ({\x},{\c/\x});
 }
\foreach \a in {-1.0,-.5,0,.5,1.0}
 \draw[ForestGreen!65!black,domain=.45:2.05,samples=50]
   plot ({sqrt(2)*\x*exp(\a)/2},{sqrt(2)*\x*exp(-\a)/2});
\node[Purple] at (2.55,2.0){constant torque};
\node[ForestGreen!60!black] at (2.65,.62){allocation lines};
\node[align=left,draw,rounded corners,fill=MidnightBlue!6] at (6.0,.4)
 {state-dependent grid:\\axes bend with the flux map\\and acquire connection terms};
\end{tikzpicture}
\caption{A curvilinear grid that follows torque and allocation structure.
Under saturation the families deform and must be tabulated or evaluated from
the local flux Jacobian.}
\label{fig:curvilinear-grid}
\end{figure}
